% META: {"id":"1SPE-PROBCOND-CO-009","chapitre":"1SPE-PROBA-COND","type_objet":"corrige","exercice_ref":"1SPE-PROBCOND-EX-009","capacites_codes":["C1"],"sources_inspiration":[],"status":"approved","origin":"generated","status_history":["generated","audited","accepted","approved"],"release_acceptance":"RELEASE_OWNER_BATCH_ACCEPTANCE","acceptance_closure_digest":"sha256:9b3ccf9a81c5520fbb7e03b7b2d3e7bf2057a02834b6d908bf3be5f49f3b553f","acceptance_receipt_digest":"sha256:4fad86f0282e0fa4e0419cca1ad6379ae7af5a280c04a4a90880f90a1c044242"}
% BEGIN-VERIFY
% from sympy import Rational, Symbol, solve
% p = Symbol('p')
% eq = 2*p/3 - Rational(1,4)
% sol = solve(eq, p)
% assert sol[0] == Rational(3, 8)
% END-VERIFY
\begin{corrige}{1SPE-PROBCOND-EX-009}

\textbf{1.} $P(A \cap B) = P(A) \times P_A(B) = p \times \dfrac{1}{3} = \dfrac{p}{3}$.

\medskip
\textbf{2.} $P_B(A) = \dfrac{P(A \cap B)}{P(B)} = \dfrac{\frac{p}{3}}{\frac{1}{2}} = \dfrac{2p}{3}$.

On a $P_B(A) = \dfrac{1}{4}$, donc $\dfrac{2p}{3} = \dfrac{1}{4}$, soit $p = \dfrac{3}{8}$.

\medskip
\textbf{3.} Vérification : $P(A \cap B) = \dfrac{3/8}{3} = \dfrac{1}{8}$. $P_A(B) = \dfrac{1/8}{3/8} = \dfrac{1}{3}$ \checkmark. $P_B(A) = \dfrac{1/8}{1/2} = \dfrac{1}{4}$ \checkmark.

\end{corrige}
